\section{Limitations and Future Work}
\label{sec:limitation}

\textbf{Cost.} The GPT-4 API incurs significant costs. It is $15 \times$ more expensive than GPT-3.5. Nevertheless, \voyager requires the quantum leap in code generation quality from GPT-4 (Fig.~\ref{fig:ablation}), which GPT-3.5 and open-source LLMs cannot provide~\cite{touvron2023llama}. 

\textbf{Inaccuracies.} Despite the iterative prompting mechanism, there are still cases where the agent gets stuck and fails to generate the correct skill. The automatic curriculum has the flexibility to reattempt this task at a later time.
Occasionally, self-verification module may also fail, such as not recognizing spider string as a success signal of beating a spider.

\textbf{Hallucinations.}
The automatic curriculum occasionally proposes unachievable tasks. For example, it may ask the agent to craft a ``copper sword" or ``copper chestplate", which are items that do not exist within the game. Hallucinations also occur during the code generation process. For instance, GPT-4 tends to use cobblestone as a fuel input, despite being an invalid fuel source in the game. Additionally, it may call functions absent in the provided control primitive APIs, leading to code execution errors.

We are confident that improvements in the GPT API models as well as novel techniques for finetuning open-source LLMs will overcome these limitations in the future.